\scanpage{23}
Mais alors ce système n'est autre que $\bar S'=\bar S'(N',e',(y'_\varepsilon)_{0\leq\varepsilon\leq e'})$. Notons que $p^e\geq N+1=(N'+1)p$ implique $p^{e-1}\geq N'+1$ i.e. $p^{e'}\geq N'+1$. Donc on est dans les $\overline{\mathrm{mes}}$ conditions que précédem\textsuperscript{t}, avec cette fois-ci $N'<N$. Par hyp. de récurrence on a donc $y'_0=0$, donc $y_0=0$, cqfd. Ouf !

\uline{Théorème 1} Soient $E$ un groupe abélien, $N\in\mathbb{N}$, $(x_i)_{i\in\mathbb{N}}$ une famille d'\unsure{éléments|élément} de $E$, telle qu'on ait
\[
\binom{j}{i}(x_j-x_i)=0\qquad\text{pour}\quad j>i\geq N
\]
Alors on a $\quad x_i=x_N\qquad\forall i\geq N$.

\uline{Dém} Quitte à localiser en tous les nbs premiers $p$, OPS $E$ un module sur \inserted{un} $\mathbb{Z}_p$. Il suffit de prouver $x_{N+1}=x_N$ (car par récurrence on \unsure{déduira|aura} de \unsure{m} $x_{N+1}=x_{N+2}$ etc). Choisissons $e\inserted{\in\mathbb{N}}$ tel que $p^e\geq N+1$, ne retenons que les $x_i$ pour $i$ de la forme $N$ ou $N+p^\varepsilon$, $0\leq\varepsilon\leq e$. Le lemme 2 nous montre que $x_{N+1}=x_N$, cqfd.

\uline{Cor 1} \unsure{Soient|Soit} $k$ un anneau commutatif, $S=k\{T\}$ l'alg. des polynômes à puiss. div., \deleted{$M$ un $k$-module\unsure{, $S$, $S^{+(N)}$}} \deleted{\unsure{$M$}} $S^{+(N)}$ l'idéal à p. div. \deleted{engendré par $T$ \unsure{soit|$i\geq N$}} défini par
\[
S^{+(N)}=\sum_{i\geq N}k.T^{(i)}.
\]
\deleted{et} Considérons, pour \unsure{tt|le|un} \deleted{\unsure{$M$}} $k$-module $M$, le $S$-module gradué
\[
\deleted{\unsure{=}}\ \unsure{w|\omega}_N(M)=M\otimes_S S^{+(N)}=\sum_{i\geq N}M.T^{(i)}
\]
Alors :
% archive sheet number: 22

